\section{Justification and Postmortem States}

\begin{quotex}
Fear God, and keep his commandments: for this is the whole duty of man. For God shall bring every work into judgment,
with every secret thing, whether it be good, or whether it be evil. \flright{Ecclesiastes 12:13-14}

In the \emph{Paradisio} of Dante's \emph{Divine Comedy}, that great mystical work, the coming of the soul
into the Presence of the Divine Light is a gradual process. One cannot come into the Divine fullness all at once; it is
too intense to be perceived. Preparation is needed. It has to be a slow process of knowing, a seeing of more and more,
a progressive revelation to the soul of God's ineffable Light. The soul then comes to a point where she
sees everything in a flash, but the distance to this point, not in spatial terms but spiritual, is great. In this
understanding, Dante is following his own experience. \flright{\textsc{Wayne Teasdale}, \emph{Essays in Mysticism}}

The Lord Jesus Christ … will judge the living and the dead \flright{2 Tim 4:1}

Man must recognize his complicity in the act of evil. \flright{\textsc{Carl Jung}, \emph{The Red Book}} 

\end{quotex}

\paragraph{Die before you Die}
The Initiate is one who strives to become a jivanmukta, someone liberated while still in the body. Or in other words,
the goal is the Beatific Vision while still in the body. That requires the three-fold Christian path of purgation,
illumination, and union. The Saint and the Initiate therefore die to this life in order to experience a spiritual
birth.

Christianity is the not religion of Jesus, but rather the religion of Christ. The modernists prefer to focus entirely on
the “historical Jesus”, the end result of which is the religion of this-worldly humanitarianism. But Jesus is the first
coming of the Logos. St. Augustine's realization that this Logos was tied also to Platonism led to his
conversion. The exoteric teaching is that there will be future comings of Christ.

However, the second coming is not just an event in linear time, but in eternity. Hence, the second coming of Christ is
the birth of the Logos in the soul. When the soul is still and purified, the Logos is born as the rebirth of the soul.

\textbf{John Ruysbroeck}, among others, also writes about the third coming of Christ, this time as Judge. The point is
not to condemn, but rather to offer healing through forgiveness.

\paragraph{General Judgment}
The General Judgment, or Last Judgment, occurs at the end of the world. At that point, all humanity will be judged.
Everyone will be made aware of each other's judgments. Moreover, the effects of one's acts and
omissions have consequences beyond the span of the individual life, usually for many generations. For example, a man
who beat his son in the 15th century may have set off a chain of events that persist into current times. For the
purposes of this note, nothing more needs be said about this now.

\paragraph{Justification and Self-Justification}
Justification is to make just. This is more than the mere “imputation” of justification as taught by the Reformation.
That false belief gained credence during a period of the loss of spiritual vision, which led to nominalism. So “imputed
justification” is justification in name only. What the Christ brings is real justification, that is, a change in
consciousness initiating the purgative stage of the spiritual path.

Such as we are, the human condition is committed to self-justification. This leaves one mired in subjectivity, and hence
it is ineffective as a means of justification. E.g., a man commits adultery, but tells his wife, “But it
didn't \emph{mean} anything.” Or else, he blames a “moment of weakness”, or ignorance of the facts of the
situation. We blame the circumstances of our birth, bad genes, an addictive habit. These and various other excuses for
self-justification seem limitless, although we seldom accept them from anyone else. True judgment is to see oneself
objectively. Without facing the truth, it is impossible to change and healing can't occur.

After spiritual rebirth, i.e., the birth of the Logos in the soul, one puts on the “mind of Christ”. This means our
identify is recentered on a higher level. Christ will then reveal our true motives and we will view our lives
objectively. We will understand how our acts and omissions have affected other people and events. The unrepentant will
resist that, but the repentant will welcome it.

A physician will tell you the precise state of your health, not to condemn you, but rather to cure you. He might suggest
a change in diet, an exercise regimen, surgery, or a medication. It is obvious that deceiving you about the state of
your health will not be of benefit. Nevertheless, it is up to the patient to follow the doctor's orders.

The situation is similar for sin, which is objectively harmful to the soul. Christ is the physician of the soul. When
unaware, or dimly aware, of the true motivations of our actions, these motivations are experienced as compulsions. And
“commandments” are experienced as unwelcome external checks on our unfulfilled desires. But the light of the Logos in
the soul illuminates one's actions. It is the means to self-observation which turns compulsions into free
acts.

\paragraph{The Esoteric Path}
To understand the esoteric path from Hell to Heaven, we turn to the initiates. Esoteric teachings are usually buried in
other structures, both to preserve them as well as to protect them from those who don't understand.
Obviously, the great exoteric Traditions contain an esoteric element. The exoteric rites, symbols, etc. conceal
— or reveal — a deeper content. They also can provide cover for small groups to peacefully
pursue their own yearnings (although there have been many exceptions.)

On a more practical level, other devices have been used. For example, the Tarot cards have a hermetic interpretation.
But they would most likely been lost if confined solely to small groups. Hence, their popular use as a means of
soothsaying has preserved them over the centuries. Obviously, most people would prefer to be told their future through
card layouts than to actually \emph{create} their own future.

Poetry may also be used for the same purpose. Dante belong to the \emph{Fedeli d'Amore}, an esoteric group
which required each of its members to write a poem. A great poem, like the \emph{Divine Comedy}, has been preserved and
studied over several centuries, although few will use it as a guide on the esoteric path. On this point we differ from
the “scholars”, who only see an allegory in the poem. This is unjustified, since Dante himself has given himself away.
The allegorical is merely the second level of interpretation (above the literal). Ultimately, the fourth level is the
anagogical or spiritual interpretation. The scholars want to understand it as an allegory describing the soul after
death. Obviously, Dante could not have truly intended that. What is being described are the states of the soul in this
life.

\paragraph{Heaven and Hell}
The Orthodox have maintained a deep tradition of the experience of Heaven and Hell:

\begin{quotex}
According to the saints, the “fire” that will consume sinners at the coming of the Kingdom of God is the same “fire”
that will shine with splendor in the saints. It is the “fire” of God's love; the “fire” of God Himself who
is Love. “For our God is a consuming fire” (Heb 12.29) who “dwells in unapproachable light” (1 Tim 6.16). For those who
love God and who love all creation in Him, the “consuming fire” of God will be radiant bliss and unspeakable delight.
For those who do not love God, and who do not love at all, this same “consuming fire” will be the cause of their
“weeping” and their “gnashing of teeth.” \flright{\itshape Heaven and Hell\footnote{\url{https://oca.org/orthodoxy/the-orthodox-faith/spirituality/the-kingdom-of-heaven/heaven-and-hell}}.} 

\end{quotex}
Thus Heaven and Hell are understood as opposed responses to the Divine Light of Christ as Judge. The saved will
experience the Light as glory, but it will be painful for the damned. So in a sense Hell is “empty”. God does not
“send” people to Hell, rather it is their choice, that is, it corresponds to their particular state of being in a
precise way.

The Initiate experiences Christ as the eternal Judge eternally, not as some abstract future event.  Dante, in the
\emph{Inferno}, describes the experience of Hell in all its various manifestations. The sins that lead to Hell are
based on malice, since there is never any intention to reform. Such men and women would repeat their acts over and over
again. Hence, Hell is a repetition, Nietzsche's “Eternal Return”, in which the same punishment is repeated
ad infinitum. That state is called the eternity of Hell.

Therefore, there is nothing unjust about it. If someone has no interest in God during life, why would that change after
death? If the response is that, after death, the soul knows. Unfortunately, that “knowing” — the “seeing
the light” — is experienced as painful, not as joy. As described in Nietzsche's Eternal Return,
the sinner repeats the sin over and over, with no intention to repent.

\paragraph{Purgatory}
Why purgatory? Purgatory is not proved by recourse to ancient writings, official documents, or cherry-picked Bible
passages. Instead we rely on the testimony of initiates. As per Dante, there is a slow process of revelation up the
Mountain of Purgatory. If, in this life, we see through a “glass darkly”, the sudden infusion of the Divine Light
typically requires a period of discomfort while getting accustomed to the Light. This is the stage of illumination:
perhaps there are rare exceptions, but the passage from purgation to union requires this stage. Even the Orthodox have
come to the same realization: the doctrine of the tollhouses\footnote{\url{https://www.gornahoor.net/?p=6343}} serves the same purposes.

The sinner in Hell is devoid of Love, so his primary motive is malice. He is malicious for its own sake, and the
concomitant false pleasure then keeps him in bondage. The soul in Purgatory is motivated by Love, although it is
deficient, disordered, or excessive. These are sins of weakness, concupiscence and ignorance. The sinner in purgatory
has pursued love through an evil means, or his love pursues a disordered end, or else the love for a good when
moderate, but bad when excessive.

These sins can be rectified: the deficient can see the unintended harmful effects of their choices, the weak can grow
stronger, the ignorant can learn from their mistakes. Thus Dante's \emph{Purgatorio} describes how the sins
in Purgatory can be overcome.

The pain and suffering of purgatory is not experienced in a negative way. It is more like the way an athlete experiences
pain during his training, or an ascetic during his self-discipline. Unlike the pains of Hell, the pain of purgatory
serve a purpose. They lead to deliverance, while the pains of Hell lead nowhere.

Consider:

\begin{itemize}
\item Accept responsibility for one's actions and their consequences 
\item Eschew self-justification 
\item See yourself as others see you. (See, for example, Schumacher's Four Fields of Knowledge\footnote{\url{https://www.gornahoor.net/?p=507}}) 
\item Recognize the consequences of your acts, including those unintended 
\item Engage in self-observation to develop awareness of all your motivations 
\end{itemize}

\flrightit{Posted on 2017-10-18 by Cologero}

\begin{center}* * *\end{center}

\begin{footnotesize}\begin{sffamily}

\texttt{Ann K on 2017-10-21 at 11:17 said: }

Thank you! A brilliant overview and very helpful to me, a catechumen in the Orthodox Church.


\hfill

\texttt{jonathansolvie on 2017-11-16 at 09:35 said: }

Thank you for this excellent article, sir. What could be a more needful matter for reflection? Why are we not taught
this in school? People are no longer taking the day of judgement seriously, believing it to be some fairytale, never
having been exposed to anything but the outermost exoteric surface of these teachings. 

“The General Judgment, or Last Judgment, occurs at the end of the world. At that point, all humanity will be judged.
Everyone will be made aware of each other's judgments.”

I want to believe this myself, but I lack certainty regarding it. How does it work metaphysically? A mere apprentice of
the Mysteries as I am, such a summary of the religious dogma, the what, doesn't contribute much to a deeper
understanding of the why and how — though I do accept Revelation. 

I find it straight forward to accept that there is a General Judgment at the end of the present human world (which I, in
accord with Guénon, presume is relatively close at hand, considering its rapidly worsening state: it is a matter of
centuries rather than of millennia now; but a more specific estimation than that must be withheld). I was going to say:
a General Judgment (from now on G.J.) for the majority of those individual souls who have found themselves
in the human state after the Fall into time is not hard to accept. What I have some difficulty with is the idea that
ALL who have lived would be included in this G.J. at the end of the world. What of those exceptional
men—great sages, adepts, yogins—who supposedly, according to their traditions, surpassed human
paradise itself and were, so to speak, ‘extinguished’ in God, or in any case, are no longer
tied up with the destiny of this particular human world? Can we expect to meet Shakyamuni or Shankara during the G.J.?
Church dogma would say yes, but is it not a generalized simplification to say so, according to a deeper esoteric
perspective? Can it be unofficially and informally admitted that there are exceptions to the general rule of the G.J.?
(I am not asking this because I expect to avoid it for my part.) 

In short: to whom does the G.J. apply and to whom may it possibly not apply? Would it be wrong to say that it applies to
souls who remain within the gravitational field' of the human state, in one of its
extensions’, without having attained the resurrection-body exceptionally early, as typified in
the Biblical tradition by Enoch, Elijah and the like, but no doubt having occurred more times than what is recorded in
that tradition, as exemplified by the rare Daoist adept whose body is said to disappear from the tomb, etc. And how
does the possibility affirmed in the East of transmigration' to non-human states, accepted by
Guénon and Schuon as possible in spite of their adoption of Sufi Islam, fit into this? 

I understand that the General Judgment is a mystery that we must probably await grasping in its fullness until we get to
experience it when the Hour strikes, but these questions still interest me.


\hfill

\texttt{Max on 2017-11-26 at 17:42 said: }

We do not take judgement seriously when we refuse to approach it inwardly as an objective experience at every moment.
All the emotional frustrations repeatedly encountered in the world should make us halt and reflect on the situation,
but instead we often double down on the same errors as if sheer velocity could release us from the downward pull of the
earth's field of gravity. The sense of dull repetition trapping many is a clue that the state of things is
not right. Faced with the spleen of modernity there comes a temptation to pursue ways of numbing and forgetting instead
of letting it act as a catalyst for rising to higher being. The same mental energy provoked through frustrations can be
turned around for good uses.

\end{sffamily}\end{footnotesize}